
\documentclass{article}
\usepackage{colortbl}
\usepackage{makecell}
\usepackage{multirow}
\usepackage{supertabular}

\begin{document}

\newcounter{utterance}

\centering \large Interaction Transcript for game `codenames', experiment `association\_difficult', episode 0 with mistral{-}large{-}3{-}azure.
\vspace{24pt}

{ \footnotesize  \setcounter{utterance}{1}
\setlength{\tabcolsep}{0pt}
\begin{supertabular}{c@{$\;$}|p{.15\linewidth}@{}p{.15\linewidth}p{.15\linewidth}p{.15\linewidth}p{.15\linewidth}p{.15\linewidth}}
    \# & $\;$A & \multicolumn{4}{c}{Game Master} & $\;\:$B\\
    \hline

    \theutterance \stepcounter{utterance}  
    & & & \multicolumn{2}{p{0.3\linewidth}}{
        \cellcolor[rgb]{0.9,0.9,0.9}{
            \makecell[{{p{\linewidth}}}]{
                \texttt{\tiny{[GM$|$GM]}}
                {'hidden': {'team': ['grzechotnik', 'czarny', 'okoń', 'miedź', 'piżama', 'stolarz', 'kurkuma', 'quad', 'fiołkowy'], 'innocent': ['traktor', 'koral', 'magenta', 'mosiądz', 'cukier', 'pieprz', 'elektryk'], 'opponent': ['szczupak', 'majtki', 'skuter', 'tymianek', 'strażak', 'żmija', 'sedan', 'aluminium'], 'assassin': ['biustonosz']}, 'revealed': {'team': {'team': [], 'innocent': [], 'opponent': [], 'assassin': []}, 'opponent': {'team': [], 'innocent': [], 'opponent': [], 'assassin': []}}}
            }
        }
    }
    & & \\ \\

    \theutterance \stepcounter{utterance}  
    & & \multicolumn{4}{p{0.6\linewidth}}{
        \cellcolor[rgb]{0.9,0.9,0.9}{
            \makecell[{{p{\linewidth}}}]{
                \texttt{\tiny{[P1$\langle$GM]}}
                \texttt{Podaj jedną wskazówkę składającą się z pojedynczego słowa, która odnosi się do niektórych słów twojej drużyny z poniższej listy. Możesz wybrać, do ilu i do jakich słów odnosi się twoja wskazówka.} \\
\texttt{Celem jest, aby twój partner odgadł wszystkie słowa twojej drużyny na podstawie podanych wskazówek, zanim drużyna przeciwna odkryje wszystkie swoje. Jeśli twój partner odgadnie słowo 'przeciwnika' lub 'rozpraszające', jego tura zgadywania kończy się i drużyna przeciwna wykonuje turę. Jeśli twój partner odgadnie słowo 'zabójcze', gra kończy się i twoja drużyna natychmiast przegrywa.} \\
\texttt{Im więcej słów odnosi się do twojej wskazówki, tym szybciej możesz wygrać grę, ale tym trudniejsze mogą być skojarzenia do odgadnięcia dla twojego partnera, więc wybieraj słowo wskazówki mądrze. Słowo wskazówki musi być semantycznie powiązane ze słowami docelowymi, nie może być jednym ze słów na listach ani zawierać ich części.} \\
\\ 
\texttt{Zawsze podawaj swoją wskazówkę składającą się z pojedynczego słowa i oddzieloną przecinkami listę powiązanych słów docelowych w następującym formacie i staraj się, aby twoje odpowiedzi były jak najkrótsze, nigdy nie dołączaj żadnego innego tekstu niż jest to wymagane w tym formacie:} \\
\\ 
\texttt{WSKAZÓWKA: <SŁOWO>} \\
\texttt{CELE: <CELE>} \\
\\ 
\texttt{Słowa twojej drużyny to: grzechotnik, czarny, okoń, miedź, piżama, stolarz, kurkuma, quad, fiołkowy.} \\
\texttt{Słowa drużyny przeciwnika to: szczupak, majtki, skuter, tymianek, strażak, żmija, sedan, aluminium.} \\
\texttt{Słowa rozpraszające to: traktor, koral, magenta, mosiądz, cukier, pieprz, elektryk.} \\
\texttt{Słowa zabójcze to: biustonosz.} \\
            }
        }
    }
    & \\ \\

    \theutterance \stepcounter{utterance}  
    & \multicolumn{4}{p{0.6\linewidth}}{
        \cellcolor[rgb]{0.9,0.9,0.9}{
            \makecell[{{p{\linewidth}}}]{
                \texttt{\tiny{[P1$\rangle$GM]}}
                \texttt{WSKAZÓWKA: kolor} \\
\texttt{CELE: czarny, fiołkowy, kurkuma} \\
            }
        }
    }
    & & \\ \\

    \theutterance \stepcounter{utterance}  
    & & & \multicolumn{2}{p{0.3\linewidth}}{
        \cellcolor[rgb]{0.9,0.9,0.9}{
            \makecell[{{p{\linewidth}}}]{
                \texttt{\tiny{[GM$|$GM]}}
                \texttt{kolor} \\
            }
        }
    }
    & & \\ \\

    \theutterance \stepcounter{utterance}  
    & & & \multicolumn{2}{p{0.3\linewidth}}{
        \cellcolor[rgb]{0.9,0.9,0.9}{
            \makecell[{{p{\linewidth}}}]{
                \texttt{\tiny{[GM$|$GM]}}
                ['czarny', 'fiołkowy', 'kurkuma']
            }
        }
    }
    & & \\ \\

    \theutterance \stepcounter{utterance}  
    & & & \multicolumn{2}{p{0.3\linewidth}}{
        \cellcolor[rgb]{0.9,0.9,0.9}{
            \makecell[{{p{\linewidth}}}]{
                \texttt{\tiny{[GM$|$GM]}}
                {'word': 'czarny', 'assignment': 'team'}
            }
        }
    }
    & & \\ \\

    \theutterance \stepcounter{utterance}  
    & & & \multicolumn{2}{p{0.3\linewidth}}{
        \cellcolor[rgb]{0.9,0.9,0.9}{
            \makecell[{{p{\linewidth}}}]{
                \texttt{\tiny{[GM$|$GM]}}
                {'word': 'fiołkowy', 'assignment': 'team'}
            }
        }
    }
    & & \\ \\

    \theutterance \stepcounter{utterance}  
    & & & \multicolumn{2}{p{0.3\linewidth}}{
        \cellcolor[rgb]{0.9,0.9,0.9}{
            \makecell[{{p{\linewidth}}}]{
                \texttt{\tiny{[GM$|$GM]}}
                {'word': 'kurkuma', 'assignment': 'team'}
            }
        }
    }
    & & \\ \\

    \theutterance \stepcounter{utterance}  
    & & \multicolumn{4}{p{0.6\linewidth}}{
        \cellcolor[rgb]{0.9,0.9,0.9}{
            \makecell[{{p{\linewidth}}}]{
                \texttt{\tiny{[GM$\rangle$P2]}}
                \texttt{Podaj oddzieloną przecinkami listę maksymalnie 3 słów z poniższej listy, które najlepiej odnoszą się lub są najbardziej związane ze słowem 'kolor'. Zawsze rozpoczynaj swoją listę odgadnięć od ZGADNIĘCIE:  i nie dołączaj żadnego innego tekstu w swojej odpowiedzi.} \\
\\ 
\texttt{tymianek, okoń, elektryk, koral, kurkuma, pieprz, miedź, stolarz, sedan, traktor, cukier, magenta, fiołkowy, szczupak, quad, żmija, piżama, strażak, czarny, majtki, biustonosz, mosiądz, skuter, aluminium, grzechotnik} \\
            }
        }
    }
    & \\ \\

    \theutterance \stepcounter{utterance}  
    & & & \multicolumn{4}{p{0.6\linewidth}}{
        \cellcolor[rgb]{0.9,0.9,0.9}{
            \makecell[{{p{\linewidth}}}]{
                \texttt{\tiny{[GM$\langle$P2]}}
                \texttt{ZGADNIĘCIE: magenta, fiołkowy, czarny} \\
            }
        }
    }
     \\ \\

    \theutterance \stepcounter{utterance}  
    & & & \multicolumn{2}{p{0.3\linewidth}}{
        \cellcolor[rgb]{0.9,0.9,0.9}{
            \makecell[{{p{\linewidth}}}]{
                \texttt{\tiny{[GM$|$GM]}}
                ['magenta', 'fiołkowy', 'czarny']
            }
        }
    }
    & & \\ \\

    \theutterance \stepcounter{utterance}  
    & & & \multicolumn{2}{p{0.3\linewidth}}{
        \cellcolor[rgb]{0.9,0.9,0.9}{
            \makecell[{{p{\linewidth}}}]{
                \texttt{\tiny{[GM$|$GM]}}
                {'word': 'magenta', 'assignment': 'innocent'}
            }
        }
    }
    & & \\ \\

    \theutterance \stepcounter{utterance}  
    & & & \multicolumn{2}{p{0.3\linewidth}}{
        \cellcolor[rgb]{0.9,0.9,0.9}{
            \makecell[{{p{\linewidth}}}]{
                \texttt{\tiny{[GM$|$GM]}}
                \texttt{magenta} \\
            }
        }
    }
    & & \\ \\

    \theutterance \stepcounter{utterance}  
    & & & \multicolumn{2}{p{0.3\linewidth}}{
        \cellcolor[rgb]{0.9,0.9,0.9}{
            \makecell[{{p{\linewidth}}}]{
                \texttt{\tiny{[GM$|$GM]}}
                {'word': 'sedan', 'assignment': 'opponent'}
            }
        }
    }
    & & \\ \\

    \theutterance \stepcounter{utterance}  
    & & & \multicolumn{2}{p{0.3\linewidth}}{
        \cellcolor[rgb]{0.9,0.9,0.9}{
            \makecell[{{p{\linewidth}}}]{
                \texttt{\tiny{[GM$|$GM]}}
                {'hidden': {'team': ['grzechotnik', 'czarny', 'okoń', 'miedź', 'piżama', 'stolarz', 'kurkuma', 'quad', 'fiołkowy'], 'innocent': ['traktor', 'koral', 'mosiądz', 'cukier', 'pieprz', 'elektryk'], 'opponent': ['szczupak', 'majtki', 'skuter', 'tymianek', 'strażak', 'żmija', 'aluminium'], 'assassin': ['biustonosz']}, 'revealed': {'team': {'team': [], 'innocent': ['magenta'], 'opponent': [], 'assassin': []}, 'opponent': {'team': [], 'innocent': [], 'opponent': ['sedan'], 'assassin': []}}}
            }
        }
    }
    & & \\ \\

    \theutterance \stepcounter{utterance}  
    & & \multicolumn{4}{p{0.6\linewidth}}{
        \cellcolor[rgb]{0.9,0.9,0.9}{
            \makecell[{{p{\linewidth}}}]{
                \texttt{\tiny{[P1$\langle$GM]}}
                \texttt{Słowo magenta zostało odgadnięte, ale to słowo niewinny. Tura twojego kolegi z drużyny się tam skończyła.} \\
\texttt{Teraz podaj kolejną wskazówkę odnoszącą się do niektórych z pozostałych słów twojego zespołu oraz listę powiązanych słów docelowych. Pamiętaj, aby rozpocząć wskazówkę od 'WSKAZÓWKA: ', wstawić nową linię i rozpocząć listę słów docelowych oddzielonych przecinkami od 'CELE: '. Uwaga: niektóre słowa zostały usunięte z list w porównaniu do poprzednich żądań.} \\
\\ 
\texttt{Twoje pozostałe słowa zespołu to: grzechotnik, czarny, okoń, miedź, piżama, stolarz, kurkuma, quad, fiołkowy.} \\
\texttt{Pozostałe słowa twojego przeciwnika to: szczupak, majtki, skuter, tymianek, strażak, żmija, sedan, aluminium.} \\
\texttt{Pozostałe słowa rozpraszające to: traktor, koral, mosiądz, cukier, pieprz, elektryk.} \\
\texttt{Pozostałe słowa zabójcy to: biustonosz.} \\
            }
        }
    }
    & \\ \\

    \theutterance \stepcounter{utterance}  
    & \multicolumn{4}{p{0.6\linewidth}}{
        \cellcolor[rgb]{0.9,0.9,0.9}{
            \makecell[{{p{\linewidth}}}]{
                \texttt{\tiny{[P1$\rangle$GM]}}
                \texttt{WSKAZÓWKA: narzędzie} \\
\texttt{CELE: stolarz, miedź} \\
            }
        }
    }
    & & \\ \\

    \theutterance \stepcounter{utterance}  
    & & & \multicolumn{2}{p{0.3\linewidth}}{
        \cellcolor[rgb]{0.9,0.9,0.9}{
            \makecell[{{p{\linewidth}}}]{
                \texttt{\tiny{[GM$|$GM]}}
                \texttt{narzędzie} \\
            }
        }
    }
    & & \\ \\

    \theutterance \stepcounter{utterance}  
    & & & \multicolumn{2}{p{0.3\linewidth}}{
        \cellcolor[rgb]{0.9,0.9,0.9}{
            \makecell[{{p{\linewidth}}}]{
                \texttt{\tiny{[GM$|$GM]}}
                ['stolarz', 'miedź']
            }
        }
    }
    & & \\ \\

    \theutterance \stepcounter{utterance}  
    & & & \multicolumn{2}{p{0.3\linewidth}}{
        \cellcolor[rgb]{0.9,0.9,0.9}{
            \makecell[{{p{\linewidth}}}]{
                \texttt{\tiny{[GM$|$GM]}}
                {'word': 'stolarz', 'assignment': 'team'}
            }
        }
    }
    & & \\ \\

    \theutterance \stepcounter{utterance}  
    & & & \multicolumn{2}{p{0.3\linewidth}}{
        \cellcolor[rgb]{0.9,0.9,0.9}{
            \makecell[{{p{\linewidth}}}]{
                \texttt{\tiny{[GM$|$GM]}}
                {'word': 'miedź', 'assignment': 'team'}
            }
        }
    }
    & & \\ \\

    \theutterance \stepcounter{utterance}  
    & & \multicolumn{4}{p{0.6\linewidth}}{
        \cellcolor[rgb]{0.9,0.9,0.9}{
            \makecell[{{p{\linewidth}}}]{
                \texttt{\tiny{[GM$\rangle$P2]}}
                \texttt{Słowo magenta zostało odgadnięte, ale to słowo niewinny. Twoja tura się tam skończyła.} \\
\texttt{Teraz podaj kolejną listę oddzieloną przecinkami, zawierającą co najmniej 1 i maksymalnie 2 słów z poniższej listy słów, które najlepiej odnoszą się lub są najbardziej związane ze słowem 'narzędzie'. Pamiętaj, aby rozpocząć odpowiedź od ZGADNIĘCIE: . Uwaga: niektóre słowa zostały usunięte z listy w porównaniu do poprzednich próśb.} \\
\\ 
\texttt{tymianek, okoń, elektryk, koral, kurkuma, pieprz, miedź, stolarz, traktor, cukier, fiołkowy, szczupak, quad, żmija, piżama, strażak, czarny, majtki, biustonosz, mosiądz, skuter, aluminium, grzechotnik} \\
            }
        }
    }
    & \\ \\

    \theutterance \stepcounter{utterance}  
    & & & \multicolumn{4}{p{0.6\linewidth}}{
        \cellcolor[rgb]{0.9,0.9,0.9}{
            \makecell[{{p{\linewidth}}}]{
                \texttt{\tiny{[GM$\langle$P2]}}
                \texttt{ZGADNIĘCIE: stolarz, traktor} \\
            }
        }
    }
     \\ \\

    \theutterance \stepcounter{utterance}  
    & & & \multicolumn{2}{p{0.3\linewidth}}{
        \cellcolor[rgb]{0.9,0.9,0.9}{
            \makecell[{{p{\linewidth}}}]{
                \texttt{\tiny{[GM$|$GM]}}
                ['stolarz', 'traktor']
            }
        }
    }
    & & \\ \\

    \theutterance \stepcounter{utterance}  
    & & & \multicolumn{2}{p{0.3\linewidth}}{
        \cellcolor[rgb]{0.9,0.9,0.9}{
            \makecell[{{p{\linewidth}}}]{
                \texttt{\tiny{[GM$|$GM]}}
                {'word': 'stolarz', 'assignment': 'team'}
            }
        }
    }
    & & \\ \\

    \theutterance \stepcounter{utterance}  
    & & & \multicolumn{2}{p{0.3\linewidth}}{
        \cellcolor[rgb]{0.9,0.9,0.9}{
            \makecell[{{p{\linewidth}}}]{
                \texttt{\tiny{[GM$|$GM]}}
                {'word': 'stolarz', 'assignment': 'team'}
            }
        }
    }
    & & \\ \\

    \theutterance \stepcounter{utterance}  
    & & & \multicolumn{2}{p{0.3\linewidth}}{
        \cellcolor[rgb]{0.9,0.9,0.9}{
            \makecell[{{p{\linewidth}}}]{
                \texttt{\tiny{[GM$|$GM]}}
                {'word': 'traktor', 'assignment': 'innocent'}
            }
        }
    }
    & & \\ \\

    \theutterance \stepcounter{utterance}  
    & & & \multicolumn{2}{p{0.3\linewidth}}{
        \cellcolor[rgb]{0.9,0.9,0.9}{
            \makecell[{{p{\linewidth}}}]{
                \texttt{\tiny{[GM$|$GM]}}
                \texttt{traktor} \\
            }
        }
    }
    & & \\ \\

    \theutterance \stepcounter{utterance}  
    & & & \multicolumn{2}{p{0.3\linewidth}}{
        \cellcolor[rgb]{0.9,0.9,0.9}{
            \makecell[{{p{\linewidth}}}]{
                \texttt{\tiny{[GM$|$GM]}}
                {'word': 'szczupak', 'assignment': 'opponent'}
            }
        }
    }
    & & \\ \\

    \theutterance \stepcounter{utterance}  
    & & & \multicolumn{2}{p{0.3\linewidth}}{
        \cellcolor[rgb]{0.9,0.9,0.9}{
            \makecell[{{p{\linewidth}}}]{
                \texttt{\tiny{[GM$|$GM]}}
                {'hidden': {'team': ['grzechotnik', 'czarny', 'okoń', 'miedź', 'piżama', 'kurkuma', 'quad', 'fiołkowy'], 'innocent': ['koral', 'mosiądz', 'cukier', 'pieprz', 'elektryk'], 'opponent': ['majtki', 'skuter', 'tymianek', 'strażak', 'żmija', 'aluminium'], 'assassin': ['biustonosz']}, 'revealed': {'team': {'team': ['stolarz'], 'innocent': ['magenta', 'traktor'], 'opponent': [], 'assassin': []}, 'opponent': {'team': [], 'innocent': [], 'opponent': ['sedan', 'szczupak'], 'assassin': []}}}
            }
        }
    }
    & & \\ \\

    \theutterance \stepcounter{utterance}  
    & & \multicolumn{4}{p{0.6\linewidth}}{
        \cellcolor[rgb]{0.9,0.9,0.9}{
            \makecell[{{p{\linewidth}}}]{
                \texttt{\tiny{[P1$\langle$GM]}}
                \texttt{Słowo stolarz zostało odgadnięte prawidłowo. Słowo traktor zostało odgadnięte, ale to słowo niewinny. Tura twojego kolegi z drużyny się tam skończyła.} \\
\texttt{Teraz podaj kolejną wskazówkę odnoszącą się do niektórych z pozostałych słów twojego zespołu oraz listę powiązanych słów docelowych. Pamiętaj, aby rozpocząć wskazówkę od 'WSKAZÓWKA: ', wstawić nową linię i rozpocząć listę słów docelowych oddzielonych przecinkami od 'CELE: '. Uwaga: niektóre słowa zostały usunięte z list w porównaniu do poprzednich żądań.} \\
\\ 
\texttt{Twoje pozostałe słowa zespołu to: grzechotnik, czarny, okoń, miedź, piżama, kurkuma, quad, fiołkowy.} \\
\texttt{Pozostałe słowa twojego przeciwnika to: szczupak, majtki, skuter, tymianek, strażak, żmija, aluminium.} \\
\texttt{Pozostałe słowa rozpraszające to: koral, mosiądz, cukier, pieprz, elektryk.} \\
\texttt{Pozostałe słowa zabójcy to: biustonosz.} \\
            }
        }
    }
    & \\ \\

    \theutterance \stepcounter{utterance}  
    & \multicolumn{4}{p{0.6\linewidth}}{
        \cellcolor[rgb]{0.9,0.9,0.9}{
            \makecell[{{p{\linewidth}}}]{
                \texttt{\tiny{[P1$\rangle$GM]}}
                \texttt{WSKAZÓWKA: pojazd} \\
\texttt{CELE: quad} \\
            }
        }
    }
    & & \\ \\

    \theutterance \stepcounter{utterance}  
    & & & \multicolumn{2}{p{0.3\linewidth}}{
        \cellcolor[rgb]{0.9,0.9,0.9}{
            \makecell[{{p{\linewidth}}}]{
                \texttt{\tiny{[GM$|$GM]}}
                \texttt{pojazd} \\
            }
        }
    }
    & & \\ \\

    \theutterance \stepcounter{utterance}  
    & & & \multicolumn{2}{p{0.3\linewidth}}{
        \cellcolor[rgb]{0.9,0.9,0.9}{
            \makecell[{{p{\linewidth}}}]{
                \texttt{\tiny{[GM$|$GM]}}
                ['quad']
            }
        }
    }
    & & \\ \\

    \theutterance \stepcounter{utterance}  
    & & & \multicolumn{2}{p{0.3\linewidth}}{
        \cellcolor[rgb]{0.9,0.9,0.9}{
            \makecell[{{p{\linewidth}}}]{
                \texttt{\tiny{[GM$|$GM]}}
                {'word': 'quad', 'assignment': 'team'}
            }
        }
    }
    & & \\ \\

    \theutterance \stepcounter{utterance}  
    & & \multicolumn{4}{p{0.6\linewidth}}{
        \cellcolor[rgb]{0.9,0.9,0.9}{
            \makecell[{{p{\linewidth}}}]{
                \texttt{\tiny{[GM$\rangle$P2]}}
                \texttt{Słowo stolarz zostało odgadnięte prawidłowo. Słowo traktor zostało odgadnięte, ale to słowo niewinny. Twoja tura się tam skończyła.} \\
\texttt{Teraz podaj kolejną listę oddzieloną przecinkami, zawierającą co najmniej 1 i maksymalnie 1 słów z poniższej listy słów, które najlepiej odnoszą się lub są najbardziej związane ze słowem 'pojazd'. Pamiętaj, aby rozpocząć odpowiedź od ZGADNIĘCIE: . Uwaga: niektóre słowa zostały usunięte z listy w porównaniu do poprzednich próśb.} \\
\\ 
\texttt{tymianek, okoń, elektryk, koral, kurkuma, pieprz, miedź, cukier, fiołkowy, quad, żmija, piżama, strażak, czarny, majtki, biustonosz, mosiądz, skuter, aluminium, grzechotnik} \\
            }
        }
    }
    & \\ \\

    \theutterance \stepcounter{utterance}  
    & & & \multicolumn{4}{p{0.6\linewidth}}{
        \cellcolor[rgb]{0.9,0.9,0.9}{
            \makecell[{{p{\linewidth}}}]{
                \texttt{\tiny{[GM$\langle$P2]}}
                \texttt{ZGADNIĘCIE: skuter} \\
            }
        }
    }
     \\ \\

    \theutterance \stepcounter{utterance}  
    & & & \multicolumn{2}{p{0.3\linewidth}}{
        \cellcolor[rgb]{0.9,0.9,0.9}{
            \makecell[{{p{\linewidth}}}]{
                \texttt{\tiny{[GM$|$GM]}}
                ['skuter']
            }
        }
    }
    & & \\ \\

    \theutterance \stepcounter{utterance}  
    & & & \multicolumn{2}{p{0.3\linewidth}}{
        \cellcolor[rgb]{0.9,0.9,0.9}{
            \makecell[{{p{\linewidth}}}]{
                \texttt{\tiny{[GM$|$GM]}}
                {'word': 'skuter', 'assignment': 'opponent'}
            }
        }
    }
    & & \\ \\

    \theutterance \stepcounter{utterance}  
    & & & \multicolumn{2}{p{0.3\linewidth}}{
        \cellcolor[rgb]{0.9,0.9,0.9}{
            \makecell[{{p{\linewidth}}}]{
                \texttt{\tiny{[GM$|$GM]}}
                \texttt{skuter} \\
            }
        }
    }
    & & \\ \\

    \theutterance \stepcounter{utterance}  
    & & & \multicolumn{2}{p{0.3\linewidth}}{
        \cellcolor[rgb]{0.9,0.9,0.9}{
            \makecell[{{p{\linewidth}}}]{
                \texttt{\tiny{[GM$|$GM]}}
                {'word': 'aluminium', 'assignment': 'opponent'}
            }
        }
    }
    & & \\ \\

    \theutterance \stepcounter{utterance}  
    & & & \multicolumn{2}{p{0.3\linewidth}}{
        \cellcolor[rgb]{0.9,0.9,0.9}{
            \makecell[{{p{\linewidth}}}]{
                \texttt{\tiny{[GM$|$GM]}}
                {'hidden': {'team': ['grzechotnik', 'czarny', 'okoń', 'miedź', 'piżama', 'kurkuma', 'quad', 'fiołkowy'], 'innocent': ['koral', 'mosiądz', 'cukier', 'pieprz', 'elektryk'], 'opponent': ['majtki', 'tymianek', 'strażak', 'żmija'], 'assassin': ['biustonosz']}, 'revealed': {'team': {'team': ['stolarz'], 'innocent': ['magenta', 'traktor'], 'opponent': ['skuter'], 'assassin': []}, 'opponent': {'team': [], 'innocent': [], 'opponent': ['sedan', 'szczupak', 'aluminium'], 'assassin': []}}}
            }
        }
    }
    & & \\ \\

    \theutterance \stepcounter{utterance}  
    & & \multicolumn{4}{p{0.6\linewidth}}{
        \cellcolor[rgb]{0.9,0.9,0.9}{
            \makecell[{{p{\linewidth}}}]{
                \texttt{\tiny{[P1$\langle$GM]}}
                \texttt{Słowo skuter zostało odgadnięte, ale to słowo przeciwnik. Tura twojego kolegi z drużyny się tam skończyła.} \\
\texttt{Teraz podaj kolejną wskazówkę odnoszącą się do niektórych z pozostałych słów twojego zespołu oraz listę powiązanych słów docelowych. Pamiętaj, aby rozpocząć wskazówkę od 'WSKAZÓWKA: ', wstawić nową linię i rozpocząć listę słów docelowych oddzielonych przecinkami od 'CELE: '. Uwaga: niektóre słowa zostały usunięte z list w porównaniu do poprzednich żądań.} \\
\\ 
\texttt{Twoje pozostałe słowa zespołu to: grzechotnik, czarny, okoń, miedź, piżama, kurkuma, quad, fiołkowy.} \\
\texttt{Pozostałe słowa twojego przeciwnika to: majtki, tymianek, strażak, żmija, aluminium.} \\
\texttt{Pozostałe słowa rozpraszające to: koral, mosiądz, cukier, pieprz, elektryk.} \\
\texttt{Pozostałe słowa zabójcy to: biustonosz.} \\
            }
        }
    }
    & \\ \\

    \theutterance \stepcounter{utterance}  
    & \multicolumn{4}{p{0.6\linewidth}}{
        \cellcolor[rgb]{0.9,0.9,0.9}{
            \makecell[{{p{\linewidth}}}]{
                \texttt{\tiny{[P1$\rangle$GM]}}
                \texttt{WSKAZÓWKA: zwierzę} \\
\texttt{CELE: grzechotnik, okoń} \\
            }
        }
    }
    & & \\ \\

    \theutterance \stepcounter{utterance}  
    & & & \multicolumn{2}{p{0.3\linewidth}}{
        \cellcolor[rgb]{0.9,0.9,0.9}{
            \makecell[{{p{\linewidth}}}]{
                \texttt{\tiny{[GM$|$GM]}}
                \texttt{zwierzę} \\
            }
        }
    }
    & & \\ \\

    \theutterance \stepcounter{utterance}  
    & & & \multicolumn{2}{p{0.3\linewidth}}{
        \cellcolor[rgb]{0.9,0.9,0.9}{
            \makecell[{{p{\linewidth}}}]{
                \texttt{\tiny{[GM$|$GM]}}
                ['grzechotnik', 'okoń']
            }
        }
    }
    & & \\ \\

    \theutterance \stepcounter{utterance}  
    & & & \multicolumn{2}{p{0.3\linewidth}}{
        \cellcolor[rgb]{0.9,0.9,0.9}{
            \makecell[{{p{\linewidth}}}]{
                \texttt{\tiny{[GM$|$GM]}}
                {'word': 'grzechotnik', 'assignment': 'team'}
            }
        }
    }
    & & \\ \\

    \theutterance \stepcounter{utterance}  
    & & & \multicolumn{2}{p{0.3\linewidth}}{
        \cellcolor[rgb]{0.9,0.9,0.9}{
            \makecell[{{p{\linewidth}}}]{
                \texttt{\tiny{[GM$|$GM]}}
                {'word': 'okoń', 'assignment': 'team'}
            }
        }
    }
    & & \\ \\

    \theutterance \stepcounter{utterance}  
    & & \multicolumn{4}{p{0.6\linewidth}}{
        \cellcolor[rgb]{0.9,0.9,0.9}{
            \makecell[{{p{\linewidth}}}]{
                \texttt{\tiny{[GM$\rangle$P2]}}
                \texttt{Słowo skuter zostało odgadnięte, ale to słowo przeciwnik. Twoja tura się tam skończyła.} \\
\texttt{Teraz podaj kolejną listę oddzieloną przecinkami, zawierającą co najmniej 1 i maksymalnie 2 słów z poniższej listy słów, które najlepiej odnoszą się lub są najbardziej związane ze słowem 'zwierzę'. Pamiętaj, aby rozpocząć odpowiedź od ZGADNIĘCIE: . Uwaga: niektóre słowa zostały usunięte z listy w porównaniu do poprzednich próśb.} \\
\\ 
\texttt{tymianek, okoń, elektryk, koral, kurkuma, pieprz, miedź, cukier, fiołkowy, quad, żmija, piżama, strażak, czarny, majtki, biustonosz, mosiądz, grzechotnik} \\
            }
        }
    }
    & \\ \\

    \theutterance \stepcounter{utterance}  
    & & & \multicolumn{4}{p{0.6\linewidth}}{
        \cellcolor[rgb]{0.9,0.9,0.9}{
            \makecell[{{p{\linewidth}}}]{
                \texttt{\tiny{[GM$\langle$P2]}}
                \texttt{ZGADNIĘCIE: okoń, żmija} \\
            }
        }
    }
     \\ \\

    \theutterance \stepcounter{utterance}  
    & & & \multicolumn{2}{p{0.3\linewidth}}{
        \cellcolor[rgb]{0.9,0.9,0.9}{
            \makecell[{{p{\linewidth}}}]{
                \texttt{\tiny{[GM$|$GM]}}
                ['okoń', 'żmija']
            }
        }
    }
    & & \\ \\

    \theutterance \stepcounter{utterance}  
    & & & \multicolumn{2}{p{0.3\linewidth}}{
        \cellcolor[rgb]{0.9,0.9,0.9}{
            \makecell[{{p{\linewidth}}}]{
                \texttt{\tiny{[GM$|$GM]}}
                {'word': 'okoń', 'assignment': 'team'}
            }
        }
    }
    & & \\ \\

    \theutterance \stepcounter{utterance}  
    & & & \multicolumn{2}{p{0.3\linewidth}}{
        \cellcolor[rgb]{0.9,0.9,0.9}{
            \makecell[{{p{\linewidth}}}]{
                \texttt{\tiny{[GM$|$GM]}}
                {'word': 'okoń', 'assignment': 'team'}
            }
        }
    }
    & & \\ \\

    \theutterance \stepcounter{utterance}  
    & & & \multicolumn{2}{p{0.3\linewidth}}{
        \cellcolor[rgb]{0.9,0.9,0.9}{
            \makecell[{{p{\linewidth}}}]{
                \texttt{\tiny{[GM$|$GM]}}
                {'word': 'żmija', 'assignment': 'opponent'}
            }
        }
    }
    & & \\ \\

    \theutterance \stepcounter{utterance}  
    & & & \multicolumn{2}{p{0.3\linewidth}}{
        \cellcolor[rgb]{0.9,0.9,0.9}{
            \makecell[{{p{\linewidth}}}]{
                \texttt{\tiny{[GM$|$GM]}}
                \texttt{żmija} \\
            }
        }
    }
    & & \\ \\

    \theutterance \stepcounter{utterance}  
    & & & \multicolumn{2}{p{0.3\linewidth}}{
        \cellcolor[rgb]{0.9,0.9,0.9}{
            \makecell[{{p{\linewidth}}}]{
                \texttt{\tiny{[GM$|$GM]}}
                {'word': 'tymianek', 'assignment': 'opponent'}
            }
        }
    }
    & & \\ \\

    \theutterance \stepcounter{utterance}  
    & & & \multicolumn{2}{p{0.3\linewidth}}{
        \cellcolor[rgb]{0.9,0.9,0.9}{
            \makecell[{{p{\linewidth}}}]{
                \texttt{\tiny{[GM$|$GM]}}
                {'hidden': {'team': ['grzechotnik', 'czarny', 'miedź', 'piżama', 'kurkuma', 'quad', 'fiołkowy'], 'innocent': ['koral', 'mosiądz', 'cukier', 'pieprz', 'elektryk'], 'opponent': ['majtki', 'strażak'], 'assassin': ['biustonosz']}, 'revealed': {'team': {'team': ['stolarz', 'okoń'], 'innocent': ['magenta', 'traktor'], 'opponent': ['skuter', 'żmija'], 'assassin': []}, 'opponent': {'team': [], 'innocent': [], 'opponent': ['sedan', 'szczupak', 'aluminium', 'tymianek'], 'assassin': []}}}
            }
        }
    }
    & & \\ \\

    \theutterance \stepcounter{utterance}  
    & & \multicolumn{4}{p{0.6\linewidth}}{
        \cellcolor[rgb]{0.9,0.9,0.9}{
            \makecell[{{p{\linewidth}}}]{
                \texttt{\tiny{[P1$\langle$GM]}}
                \texttt{Słowo okoń zostało odgadnięte prawidłowo. Słowo żmija zostało odgadnięte, ale to słowo przeciwnik. Tura twojego kolegi z drużyny się tam skończyła.} \\
\texttt{Teraz podaj kolejną wskazówkę odnoszącą się do niektórych z pozostałych słów twojego zespołu oraz listę powiązanych słów docelowych. Pamiętaj, aby rozpocząć wskazówkę od 'WSKAZÓWKA: ', wstawić nową linię i rozpocząć listę słów docelowych oddzielonych przecinkami od 'CELE: '. Uwaga: niektóre słowa zostały usunięte z list w porównaniu do poprzednich żądań.} \\
\\ 
\texttt{Twoje pozostałe słowa zespołu to: grzechotnik, czarny, miedź, piżama, kurkuma, quad, fiołkowy.} \\
\texttt{Pozostałe słowa twojego przeciwnika to: majtki, tymianek, strażak.} \\
\texttt{Pozostałe słowa rozpraszające to: koral, mosiądz, cukier, pieprz, elektryk.} \\
\texttt{Pozostałe słowa zabójcy to: biustonosz.} \\
            }
        }
    }
    & \\ \\

    \theutterance \stepcounter{utterance}  
    & \multicolumn{4}{p{0.6\linewidth}}{
        \cellcolor[rgb]{0.9,0.9,0.9}{
            \makecell[{{p{\linewidth}}}]{
                \texttt{\tiny{[P1$\rangle$GM]}}
                \texttt{WSKAZÓWKA: metal} \\
\texttt{CELE: miedź} \\
            }
        }
    }
    & & \\ \\

    \theutterance \stepcounter{utterance}  
    & & & \multicolumn{2}{p{0.3\linewidth}}{
        \cellcolor[rgb]{0.9,0.9,0.9}{
            \makecell[{{p{\linewidth}}}]{
                \texttt{\tiny{[GM$|$GM]}}
                \texttt{metal} \\
            }
        }
    }
    & & \\ \\

    \theutterance \stepcounter{utterance}  
    & & & \multicolumn{2}{p{0.3\linewidth}}{
        \cellcolor[rgb]{0.9,0.9,0.9}{
            \makecell[{{p{\linewidth}}}]{
                \texttt{\tiny{[GM$|$GM]}}
                ['miedź']
            }
        }
    }
    & & \\ \\

    \theutterance \stepcounter{utterance}  
    & & & \multicolumn{2}{p{0.3\linewidth}}{
        \cellcolor[rgb]{0.9,0.9,0.9}{
            \makecell[{{p{\linewidth}}}]{
                \texttt{\tiny{[GM$|$GM]}}
                {'word': 'miedź', 'assignment': 'team'}
            }
        }
    }
    & & \\ \\

    \theutterance \stepcounter{utterance}  
    & & \multicolumn{4}{p{0.6\linewidth}}{
        \cellcolor[rgb]{0.9,0.9,0.9}{
            \makecell[{{p{\linewidth}}}]{
                \texttt{\tiny{[GM$\rangle$P2]}}
                \texttt{Słowo okoń zostało odgadnięte prawidłowo. Słowo żmija zostało odgadnięte, ale to słowo przeciwnik. Twoja tura się tam skończyła.} \\
\texttt{Teraz podaj kolejną listę oddzieloną przecinkami, zawierającą co najmniej 1 i maksymalnie 1 słów z poniższej listy słów, które najlepiej odnoszą się lub są najbardziej związane ze słowem 'metal'. Pamiętaj, aby rozpocząć odpowiedź od ZGADNIĘCIE: . Uwaga: niektóre słowa zostały usunięte z listy w porównaniu do poprzednich próśb.} \\
\\ 
\texttt{elektryk, koral, kurkuma, pieprz, miedź, cukier, fiołkowy, quad, piżama, strażak, czarny, majtki, biustonosz, mosiądz, grzechotnik} \\
            }
        }
    }
    & \\ \\

    \theutterance \stepcounter{utterance}  
    & & & \multicolumn{4}{p{0.6\linewidth}}{
        \cellcolor[rgb]{0.9,0.9,0.9}{
            \makecell[{{p{\linewidth}}}]{
                \texttt{\tiny{[GM$\langle$P2]}}
                \texttt{ZGADNIĘCIE: miedź} \\
            }
        }
    }
     \\ \\

    \theutterance \stepcounter{utterance}  
    & & & \multicolumn{2}{p{0.3\linewidth}}{
        \cellcolor[rgb]{0.9,0.9,0.9}{
            \makecell[{{p{\linewidth}}}]{
                \texttt{\tiny{[GM$|$GM]}}
                ['miedź']
            }
        }
    }
    & & \\ \\

    \theutterance \stepcounter{utterance}  
    & & & \multicolumn{2}{p{0.3\linewidth}}{
        \cellcolor[rgb]{0.9,0.9,0.9}{
            \makecell[{{p{\linewidth}}}]{
                \texttt{\tiny{[GM$|$GM]}}
                {'word': 'miedź', 'assignment': 'team'}
            }
        }
    }
    & & \\ \\

    \theutterance \stepcounter{utterance}  
    & & & \multicolumn{2}{p{0.3\linewidth}}{
        \cellcolor[rgb]{0.9,0.9,0.9}{
            \makecell[{{p{\linewidth}}}]{
                \texttt{\tiny{[GM$|$GM]}}
                {'word': 'miedź', 'assignment': 'team'}
            }
        }
    }
    & & \\ \\

    \theutterance \stepcounter{utterance}  
    & & & \multicolumn{2}{p{0.3\linewidth}}{
        \cellcolor[rgb]{0.9,0.9,0.9}{
            \makecell[{{p{\linewidth}}}]{
                \texttt{\tiny{[GM$|$GM]}}
                {'word': 'strażak', 'assignment': 'opponent'}
            }
        }
    }
    & & \\ \\

    \theutterance \stepcounter{utterance}  
    & & & \multicolumn{2}{p{0.3\linewidth}}{
        \cellcolor[rgb]{0.9,0.9,0.9}{
            \makecell[{{p{\linewidth}}}]{
                \texttt{\tiny{[GM$|$GM]}}
                {'hidden': {'team': ['grzechotnik', 'czarny', 'piżama', 'kurkuma', 'quad', 'fiołkowy'], 'innocent': ['koral', 'mosiądz', 'cukier', 'pieprz', 'elektryk'], 'opponent': ['majtki'], 'assassin': ['biustonosz']}, 'revealed': {'team': {'team': ['stolarz', 'okoń', 'miedź'], 'innocent': ['magenta', 'traktor'], 'opponent': ['skuter', 'żmija'], 'assassin': []}, 'opponent': {'team': [], 'innocent': [], 'opponent': ['sedan', 'szczupak', 'aluminium', 'tymianek', 'strażak'], 'assassin': []}}}
            }
        }
    }
    & & \\ \\

    \theutterance \stepcounter{utterance}  
    & & \multicolumn{4}{p{0.6\linewidth}}{
        \cellcolor[rgb]{0.9,0.9,0.9}{
            \makecell[{{p{\linewidth}}}]{
                \texttt{\tiny{[P1$\langle$GM]}}
                \texttt{Słowo miedź zostało odgadnięte prawidłowo. Tura twojego kolegi z drużyny się tam skończyła.} \\
\texttt{Teraz podaj kolejną wskazówkę odnoszącą się do niektórych z pozostałych słów twojego zespołu oraz listę powiązanych słów docelowych. Pamiętaj, aby rozpocząć wskazówkę od 'WSKAZÓWKA: ', wstawić nową linię i rozpocząć listę słów docelowych oddzielonych przecinkami od 'CELE: '. Uwaga: niektóre słowa zostały usunięte z list w porównaniu do poprzednich żądań.} \\
\\ 
\texttt{Twoje pozostałe słowa zespołu to: grzechotnik, czarny, piżama, kurkuma, quad, fiołkowy.} \\
\texttt{Pozostałe słowa twojego przeciwnika to: majtki, strażak.} \\
\texttt{Pozostałe słowa rozpraszające to: koral, mosiądz, cukier, pieprz, elektryk.} \\
\texttt{Pozostałe słowa zabójcy to: biustonosz.} \\
            }
        }
    }
    & \\ \\

    \theutterance \stepcounter{utterance}  
    & \multicolumn{4}{p{0.6\linewidth}}{
        \cellcolor[rgb]{0.9,0.9,0.9}{
            \makecell[{{p{\linewidth}}}]{
                \texttt{\tiny{[P1$\rangle$GM]}}
                \texttt{WSKAZÓWKA: ubranie} \\
\texttt{CELE: piżama, fiołkowy} \\
            }
        }
    }
    & & \\ \\

    \theutterance \stepcounter{utterance}  
    & & & \multicolumn{2}{p{0.3\linewidth}}{
        \cellcolor[rgb]{0.9,0.9,0.9}{
            \makecell[{{p{\linewidth}}}]{
                \texttt{\tiny{[GM$|$GM]}}
                \texttt{ubranie} \\
            }
        }
    }
    & & \\ \\

    \theutterance \stepcounter{utterance}  
    & & & \multicolumn{2}{p{0.3\linewidth}}{
        \cellcolor[rgb]{0.9,0.9,0.9}{
            \makecell[{{p{\linewidth}}}]{
                \texttt{\tiny{[GM$|$GM]}}
                ['piżama', 'fiołkowy']
            }
        }
    }
    & & \\ \\

    \theutterance \stepcounter{utterance}  
    & & & \multicolumn{2}{p{0.3\linewidth}}{
        \cellcolor[rgb]{0.9,0.9,0.9}{
            \makecell[{{p{\linewidth}}}]{
                \texttt{\tiny{[GM$|$GM]}}
                {'word': 'piżama', 'assignment': 'team'}
            }
        }
    }
    & & \\ \\

    \theutterance \stepcounter{utterance}  
    & & & \multicolumn{2}{p{0.3\linewidth}}{
        \cellcolor[rgb]{0.9,0.9,0.9}{
            \makecell[{{p{\linewidth}}}]{
                \texttt{\tiny{[GM$|$GM]}}
                {'word': 'fiołkowy', 'assignment': 'team'}
            }
        }
    }
    & & \\ \\

    \theutterance \stepcounter{utterance}  
    & & \multicolumn{4}{p{0.6\linewidth}}{
        \cellcolor[rgb]{0.9,0.9,0.9}{
            \makecell[{{p{\linewidth}}}]{
                \texttt{\tiny{[GM$\rangle$P2]}}
                \texttt{Słowo miedź zostało odgadnięte prawidłowo. Twoja tura się tam skończyła.} \\
\texttt{Teraz podaj kolejną listę oddzieloną przecinkami, zawierającą co najmniej 1 i maksymalnie 2 słów z poniższej listy słów, które najlepiej odnoszą się lub są najbardziej związane ze słowem 'ubranie'. Pamiętaj, aby rozpocząć odpowiedź od ZGADNIĘCIE: . Uwaga: niektóre słowa zostały usunięte z listy w porównaniu do poprzednich próśb.} \\
\\ 
\texttt{elektryk, koral, kurkuma, pieprz, cukier, fiołkowy, quad, piżama, czarny, majtki, biustonosz, mosiądz, grzechotnik} \\
            }
        }
    }
    & \\ \\

    \theutterance \stepcounter{utterance}  
    & & & \multicolumn{4}{p{0.6\linewidth}}{
        \cellcolor[rgb]{0.9,0.9,0.9}{
            \makecell[{{p{\linewidth}}}]{
                \texttt{\tiny{[GM$\langle$P2]}}
                \texttt{ZGADNIĘCIE: piżama, majtki} \\
            }
        }
    }
     \\ \\

    \theutterance \stepcounter{utterance}  
    & & & \multicolumn{2}{p{0.3\linewidth}}{
        \cellcolor[rgb]{0.9,0.9,0.9}{
            \makecell[{{p{\linewidth}}}]{
                \texttt{\tiny{[GM$|$GM]}}
                ['piżama', 'majtki']
            }
        }
    }
    & & \\ \\

    \theutterance \stepcounter{utterance}  
    & & & \multicolumn{2}{p{0.3\linewidth}}{
        \cellcolor[rgb]{0.9,0.9,0.9}{
            \makecell[{{p{\linewidth}}}]{
                \texttt{\tiny{[GM$|$GM]}}
                {'word': 'piżama', 'assignment': 'team'}
            }
        }
    }
    & & \\ \\

    \theutterance \stepcounter{utterance}  
    & & & \multicolumn{2}{p{0.3\linewidth}}{
        \cellcolor[rgb]{0.9,0.9,0.9}{
            \makecell[{{p{\linewidth}}}]{
                \texttt{\tiny{[GM$|$GM]}}
                {'word': 'piżama', 'assignment': 'team'}
            }
        }
    }
    & & \\ \\

    \theutterance \stepcounter{utterance}  
    & & & \multicolumn{2}{p{0.3\linewidth}}{
        \cellcolor[rgb]{0.9,0.9,0.9}{
            \makecell[{{p{\linewidth}}}]{
                \texttt{\tiny{[GM$|$GM]}}
                {'word': 'majtki', 'assignment': 'opponent'}
            }
        }
    }
    & & \\ \\

    \theutterance \stepcounter{utterance}  
    & & & \multicolumn{2}{p{0.3\linewidth}}{
        \cellcolor[rgb]{0.9,0.9,0.9}{
            \makecell[{{p{\linewidth}}}]{
                \texttt{\tiny{[GM$|$GM]}}
                \texttt{majtki} \\
            }
        }
    }
    & & \\ \\

    \theutterance \stepcounter{utterance}  
    & & & \multicolumn{2}{p{0.3\linewidth}}{
        \cellcolor[rgb]{0.9,0.9,0.9}{
            \makecell[{{p{\linewidth}}}]{
                \texttt{\tiny{[GM$|$GM]}}
                \texttt{opponent has won} \\
            }
        }
    }
    & & \\ \\

\end{supertabular}
}

\end{document}
